\begin{table}[htbp]
\centering
\small
\caption{Mean MFAI by industrial group, relative to city mean.}
\label{tab:s4b_group_relative}
\begin{tabular}{lrrrrrrrr}
\toprule
City & G1 Food, beverages, tobacco & G2 Textiles, apparel, leather & G3 Wood, paper, printing & G4 Petroleum, chemicals, pharma & G5 Rubber, plastics, minerals & G6 Basic and fabricated metals & G7 Electronics, machinery & G8 Transport eq., furniture, other \\
\midrule
Patna & 0.985000 & 0.996000 & 0.988000 & 1.015000 & 0.993000 & 1.013000 & 1.008000 & 1.000000 \\
Indore & 0.984000 & 1.009000 & 1.006000 & 1.001000 & 0.995000 & 1.009000 & 1.002000 & 1.004000 \\
Varanasi & 0.979000 & 0.994000 & 0.992000 & 1.024000 & 1.017000 & 1.029000 & 1.005000 & 1.003000 \\
Kochi & 0.988000 & 1.008000 & 0.987000 & 1.009000 & 0.976000 & 1.020000 & 1.033000 & 1.008000 \\
Kolkata & 0.999000 & 0.986000 & 1.004000 & 1.008000 & 1.004000 & 1.012000 & 1.009000 & 1.001000 \\
\bottomrule
\end{tabular}
\begin{tablenotes}\footnotesize
\item A value of 1.000 is the city average. All forty city–group combinations lie between 0.976 and 1.033, corroborating the 1.0\% between-sector variance share in Table 10.
\end{tablenotes}
\end{table}
